% Related Work
\section{Related Work}
\label{sec:related}

We organize related work by the detection granularity and the role of causality in the method.

\paragraph{Provenance-Based Intrusion Detection.}

Early provenance-based detectors operated at coarse granularities. UNICORN~\cite{unicorn2020} applies Weisfeiler-Lehman graph kernels and histogram sketching to classify entire provenance graphs, relying on an evolutionary model of system behavior. PROGRAPHER~\cite{yang2023prographer} partitions provenance into temporal snapshots, applies graph2vec for snapshot embedding, and uses TextRCNN for sequence-level anomaly detection. While PROGRAPHER incorporates temporal sequence modeling, it operates at the snapshot level---individual event causality within snapshots is lost.

At finer granularities, MAGIC~\cite{jia2024magic} uses a Graph Masked Auto-Encoder (GMAE) to learn benign entity representations and detects anomalies via KNN density estimation. ThreaTrace~\cite{wang2022threatrace} applies GraphSAGE for node-level anomaly detection with type-specific models. FLASH~\cite{flash2024} combines Word2Vec semantic encoding with GNN contextual encoding for node-level detection. These methods identify \emph{which entities} are suspicious, but not \emph{how they are causally connected}---the analyst receives a list of suspicious nodes, not an attack narrative.

KAIROS~\cite{cheng2024kairos} is the closest predecessor to our work in terms of input pipeline. It uses a Temporal Graph Network (TGN) to compute per-edge reconstruction loss, then applies statistical thresholding within 15-minute windows to identify anomalous time periods. KAIROS and our work address complementary problems: KAIROS produces anomalous time windows for triage, whereas our work produces causal event chains for investigation. In our current instantiation, we adopt KAIROS's pretrained TGN as the edge-level scorer that supplies seed edges to our chain extractor; our framework, however, is agnostic to this choice, and any per-edge anomaly score (e.g., from EdgeTrace~\cite{edgetrace2025} or a future specialized scorer) could play the same role.

\paragraph{Transformers on Provenance Graphs.}

A recent wave of work applies Transformer architectures to provenance data. EagleEye~\cite{eagleeye2024} linearizes provenance events into fixed-size windows and uses a BERT-tiny encoder for malware classification, demonstrating that self-attention can surface suspicious events for analyst review. Sentient~\cite{sentient2026} employs a Graph Transformer with Laplacian positional encoding and an Intent Analysis Module based on Mamba-2 for behavioral logic association. GET-AID~\cite{getaid2025} proposes a two-stage graph encoder with node-level and event-level TransformerConv attention across temporal subgraphs. PanThreat~\cite{panthreat2025} and LogShield~\cite{logshield2024} similarly apply Transformer variants to provenance or log sequences.

Our work occupies a different layer in the design space than these methods. Where the Transformer-on-provenance line treats segmentation as a preprocessing step and asks how a Transformer should consume the resulting sequences, we ask the prior question of how the sequences themselves should be defined. We posit, and demonstrate empirically through ablation, that the structure on which a Transformer attends is causal in nature, and that extracting causally-coherent chains rather than time windows, random walks, or temporal subgraphs improves chain-level detection regardless of the downstream model architecture. Our contribution is therefore complementary to advances in Transformer architectures for provenance: any of the methods above could, in principle, consume the chains we extract.

\paragraph{Attack Investigation and Path Reconstruction.}

Rule-based attack investigation systems---Holmes~\cite{holmes2019}, Poirot~\cite{poirot2019}, and SLEUTH---use expert-defined attack patterns to trace causal paths through provenance graphs. While effective for known attack patterns, they cannot generalize to novel attacks. SPARSE~\cite{sparse2024} constructs Suspicious Semantic Graphs and performs path-level analysis using semantic tags and threat intelligence---a hybrid of rule-based and learned approaches.

Recent ML-based methods perform path reconstruction as post-processing. SLOT~\cite{slot2025} uses graph reinforcement learning for anomaly detection, then applies label propagation clustering to construct attack chains. CAGE~\cite{cage2025} assigns dynamic causal weights to provenance edges via Q-learning and identifies attack paths through weighted subgraph extraction. EdgeTrace~\cite{edgetrace2025} uses a masked graph autoencoder for edge-level anomaly scoring, followed by semantic clustering for causal path inference.

In all of these methods, chain construction is a \emph{downstream step}: the chain is reconstructed after node- or edge-level detection has already produced its alerts, and chain quality is therefore bounded by upstream recall---a single missed node breaks the chain irrecoverably. Our method takes the chain itself as the detection unit, extracted upstream under a model-free coherence definition; downstream sequence modeling then operates over chains that are causally coherent by construction, allowing weak individual signals along a chain to be amplified by self-attention over the full causal context rather than discarded as isolated outliers.

\paragraph{Causal Modeling in Security.}

FALCON~\cite{falcon2025} applies federated causal learning across heterogeneous audit logs, using bidirectional Transformers for cross-domain attack investigation. While FALCON shares our interest in causal provenance, its focus is on privacy-preserving cross-domain learning rather than chain-level detection granularity. More broadly, causal inference techniques have been applied to intrusion detection~\cite{causalids}, but primarily as a means of improving detection accuracy rather than as a framework for defining and measuring chain quality.

\paragraph{Positioning summary.}

Across the four lines of work above, chain quality is either implicitly defined by the segmentation rule (time windows, random walks, temporal subgraphs), or by handcrafted attack patterns (Holmes, Poirot), or treated as a downstream artifact of node- or edge-level detection (SLOT, CAGE, EdgeTrace). What is missing is a definition of chain quality that is \emph{model-free, attack-agnostic, and computable directly on the kernel-recorded graph}---an under-explored direction we take up here. Operationally, we keep extraction (where causality is enforced) and sequence modeling (where semantic deviation is measured) at separate layers of the pipeline, so that improvements in either layer---e.g., the Transformer architectures of EagleEye, Sentient, or GET-AID, or specialized edge-level scorers---are complementary rather than competing with our contribution.
